\chapter{Ուսանող 2 — Ծրագրային ապահովում / Ալգորիթմներ}

\textbf{Անհատական թեմա.}\\
\emph{FPGA-ից ստացվող սենսորային տվյալների պահպանման, վիզուալիզացիայի և վերլուծության
ծրագրային հարթակի մշակում.}

\textbf{Հիմնական պարտականությունները.}
\begin{enumerate}[leftmargin=1.6em]
  \item Տվյալների բազայի ճարտարապետության նախագծում։
  \item Backend/API-ի մշակում և FPGA/PS տվյալների ընդունում։
  \item Փորձերի և սենսորային նմուշների պահպանում։
  \item Web հավելվածի մշակում՝ փորձի կոնֆիգուրացիայի ինտերֆեյսով։
  \item Կոնֆիգուրացիոն հրամանների ուղարկում AUP ZU3-ին։
  \item Live և պատմական տվյալների վիզուալիզացիա, փորձերի համեմատում։
  \item Վիճակագրական վերլուծության իրականացում (mean, variance, std, և այլն)։
  \item Վերլուծության/մոդելավորման բաղադրիչի իրականացում։
  \item (Ընտրովի) անոմալիաների հայտնաբերում կամ կանխատեսում (prediction)։
  \item (Ընտրովի) վերլուծության արդյունքներով FPGA-ի պարամետրերի ավտոմատ փոփոխում։
\end{enumerate}

Երկրորդ ուսանողի աշխատանքը պետք է մնա բավականաչափ անկախ RTL-ի մշակումից. հիմնական
շեշտը՝ տվյալների բազա, Backend, Web հարթակ, վիզուալիզացիա, վիճակագրություն և
վերլուծություն/մոդելավորում։ Հետաքրքիր համեմատություն է առաջանում FPGA-ի իրական
ժամանակի պարզ հայտնաբերման (threshold) և սերվերային ավելի բարդ, պատմական տվյալների
վրա հիմնված վերլուծության միջև։
